\chapter{Torticollis}

\section*{Definition}
Torticollis (wry neck) is a clinical deformity characterized by lateral flexion of the head toward the affected side with contralateral rotation, caused by abnormal contraction or shortening of the sternocleidomastoid muscle. It may be congenital or acquired, with diagnostic thresholds varying by etiology.

\section*{Pathophysiology}
Congenital muscular torticollis results from intramuscular fibrosis of the sternocleidomastoid, possibly due to intrauterine positioning, birth trauma, or compartment syndrome. Acquired torticollis involves prolonged sternocleidomastoid spasm triggered by inflammation, trauma, or neurologic irritation of the spinal accessory nerve (CN XI).

\section*{Etiology}
\begin{itemize}
  \item \textbf{Congenital:} Congenital muscular torticollis from sternocleidomastoid fibromatosis, Klippel-Feil syndrome (cervical vertebral fusion)
  \item \textbf{Traumatic:} Cervical spine fracture or dislocation, atlantoaxial rotatory subluxation, soft tissue injury
  \item \textbf{Infectious:} Retropharyngeal abscess, cervical lymphadenitis, tonsillitis (Grisel's syndrome)
  \item \textbf{Neurologic:} Dystonic reaction to medications (antipsychotics, metoclopramide), cervical dystonia (idiopathic spasmodic torticollis)
  \item \textbf{Other:} Rheumatoid arthritis affecting the cervical spine, head/neck tumors, benign paroxysmal torticollis of infancy
\end{itemize}

\section*{Indications for Evaluation}
Seek care if:
\begin{itemize}
  \item Acute onset following trauma with neck pain or neurologic deficits
  \item Fever, meningismus, or signs of infection (pharyngitis, lymphadenopathy)
  \item Persistent torticollis beyond 1 week without identifiable cause
  \item Recent initiation of antipsychotic or antiemetic medication (rule out acute dystonic reaction)
  \item Worsening head tilt in an infant with no improvement after 6 months of stretching
\end{itemize}

\section*{Related Symptoms}
\begin{itemize}
  \item Neck pain, restricted cervical range of motion, headache, shoulder elevation, facial asymmetry (infants)
\end{itemize}
\textbf{Note:} Acute torticollis with fever and odynophagia should prompt urgent evaluation for retropharyngeal abscess, particularly in children.

\section*{Self-Care / Home Management}
\begin{itemize}
  \item Gentle passive stretching of the sternocleidomastoid muscle (chin to shoulder, ear to opposite shoulder)
  \item Heat packs applied to the affected muscle for 15--20 minutes to reduce spasm
  \item NSAIDs (ibuprofen, naproxen) for pain and inflammation
  \item Cervical collar use limited to 24--48 hours for acute spasm; prolonged use may weaken muscles
  \item Prone positioning and active neck rotation exercises in infants with congenital torticollis
\end{itemize}

\section*{Diagnostic Approach}
\begin{itemize}
  \item Physical exam: assess cervical range of motion, palpate sternocleidomastoid for fibrosis or mass, evaluate for lymphadenopathy or pharyngeal erythema
  \item Cervical spine X-rays (AP, lateral, odontoid views) to exclude fracture, subluxation, or congenital anomaly
  \item CT with contrast if retropharyngeal abscess or deep neck infection is suspected
  \item MRI for persistent torticollis with neurologic signs or suspected spinal cord pathology
  \item Ultrasound of sternocleidomastoid in infants to assess for fibromatosis colli
\end{itemize}

\section*{Treatment / Management Overview}
\begin{itemize}
  \item Passive stretching and physical therapy as first-line for congenital muscular torticollis; 90\% resolve with conservative treatment by age 1
  \item Antibiotics and surgical drainage for retropharyngeal abscess
  \item Anticholinergic medications (benztropine, diphenhydramine) for acute dystonic reactions
  \item Botulinum toxin type A injections for idiopathic cervical dystonia
  \item Surgical release of the sternocleidomastoid (unipolar or bipolar) if conservative therapy fails beyond 12 months of age
\end{itemize}

\clearpage
